\chapter{State of the Art}

At the beginning of each chapter, a description should introduce the reader to the content of the chapter. The description should explain to the reader the layout of the chapter, the contribution that the chapter makes to the overall dissertation and the contribution of the individual sections towards the overall chapter.

From the perspective of this document forming part of your degree, this chapter should demonstrate to the reader your knowledge of the area of your dissertation project. It should present your knowledge in a coherent and detailed form. The reader should understand that you have in-depth knowledge of the area of the dissertation without being overloaded with information.

\section{Background}

A section on the background of the dissertation should provide the reader with an introduction to existing technologies and concepts that form the basis of the
work presented in the dissertation.

\subsection{Functional Programming}
  \href{https://www.cs.kent.ac.uk/people/staff/dat/miranda/whyfp90.pdf}{Why Functional Programming Matters} [1]
    This paper details the benefits of functional programming. Specifically it is said that they have greater capacity for modular design and that the paradigm of lazy evaluation in functional languages allow for simpler code. The reasons detailed within drove the motivation to convert from Python to Haskell  

    The reason given for why functional languages have greater capacity for modular design is through their ability to take functions as arguments to other functions.
    Side effects not being allowed in functional languages create a system where every function has a single well defined output. functions are inputted to other functions rather than have a function do multiple jobs. this allows for the easy modification of lower level functions and higher order functions do not need to worry how the lower levels work, they can be treated as black boxes. when all functions have very specific jobs, they can be reused far more often than a function performing multiple jobs, hence improving modularity.

    Lazy evaluation is a powerful principle. values are only computed when needed. infinite series can be generated and the selectors will compute only the values needed. additionally this separation of generation and selection further improves the modularity of the system. this is a very useful property of these languages for this project in particular, as for this test generation system it allows for generation of a very large number of tests without of running out of memory.
    An example of lazy evaluation being helpful is in evaluating the newton-raphson approximation of square roots. the process is an inherently infinite convergent algorithm where each approximation is better than the last. an imperative langauge would cut off the loop when two approximation differ by no more than some predefined value. Lazy evaluation however allows for the infininite series to be mapped, and then taken to any level of precision once a function calls for it.

    Hughes' arguments were influential in justifying the decision to change part of the codebase to be functional. the increased modularity would aid extending the codebase
    
\vspace{25}
\href{https://www.microsoft.com/en-us/research/wp-content/uploads/2016/07/history.pdf}{A History of Haskell} 
clearly gives a way to write about haskell and its history background and relevance/use to a project like this, more justificaiton for its use

Use it to explain why you chose Haskell:

     Strong static typing and type classes ensure your language backends are complete at compile time.
     Purity and modularity make it easier to reason about transformations from SPIN output to abstract test actions.

Haskell was created at a time where there was a dozen non-strict, functional programming languages in circulation. A conference was was held and it was decided that the lack of a common language was hampering widespread adoption of these class of languages. to remedy this, the committe decided on the creation of a language to unite the space.
the specific principles decided upon at this committee are responsible for its utility to the project at hand. The benefits talked about above still apply due to its nature as a funcitonal programming language but additionally, haskell in particular has its own advantages such as

Purity by default
    other funcitonal languages (Lisps, Scheme, Clojure), have purity as style, side effects are allowed to be used freely. 
    in haskell side effects have to programmed explicitly as functions are pure by default.

    this gives the advantages that functions are easier to reason about. just from the type signature you can  know that the function can throw exception or access the outside world by standard input output.

    it makes it easier to refactor, as when changing code you dont have to worry about affecting mutable varialbles or changing state in a subtle way.

Monads
    monads are the mechanism haskell uses to control side effects and they 


Type classes and ad hoc polymorphism
    type classes are a powerful mechanism in haskell in that they allow the programmer to define an abstract interface with mutliple concrete implementations. For this project this manifests as an abstract "test action" with concrete implementations in various different languages.

    Type classes are flexible in what fields are allowable. in addition haskell enforces completeness in type classes at compile time. this means that for all fields defined for a type class they must be implemetnted



    


\subsection{Multi-Language Systems}
    https://chapering.github.io/pubs/tse24haoran.pdf

    analyses stack overflow posts to answer what developers find hard about multingual development. I used it to learn about the methods of multingual development and solutions to challenges. it details that language interfacing is the biggest challenge faced. there is two method of interfacing, explicit and implicit. explicit interfacing involves the language invoking functions from the second language in the style of same language imports. implicit interfacing involves indirect methods of communicating eg message queues which is what I am using due to the fact that the explicit interface implementations for Haskell to python were limited as will be talked about later.

the paper also suggest more specific solutions for problems. One thing that the developers found tripped developers up a lot was data mismatches in function calls. the solution was to be explicit with all types to ensure no discrepancies    
     https://dl.acm.org/doi/abs/10.5555/3370272.3370280?download=true



     this one doesn't actually let you read it so find another one 
     two papers that talk about the benefits and challenges of multi language systems

\subsection{automated test generation using model checkers}
https://openresearch.surrey.ac.uk/esploro/outputs/conferencePresentation/A-model-checking-based-test-case/99516972702346/filesAndLinks?index=0
this is looking like a paper that does the same sorta thing that andrew butterfields thing does but like way earlier

tjis paper looks at 
Business Process Execution Language, a language that is used to describe web service composition behaviour.
due to its use of concurrency and hierarchy, it was difficult to verify manually. so the author came to a similar solution to ourselves and applied model checkers to automatically create test suites. their approach contrasts ours in theat they use NuSMV instead of Promela



How to tie it in: 

     Use it as a more general example of “model checking as test generator”, not just for your RTEMS context.
     You can say: just as they use a model checker to derive meaningful test sequences for web services, you use SPIN to generate test sequences for RTEMS, then map those to concrete test code in C/Ada/Rust.
     


          
\section{Closely-Related Work}

Work in research areas tends to address a number of specific aspects. Ideally, the discussion of published research should focus on the aspects that have been addressed by various publications - and not a discussion of the individual publications.

For example, if the topic would be a discussion of work on programming languages, the subsections of the related work could be discussions of object orientation and its realisation in various languages or the use of lambda functions by these languages.

\item[\href{https://publications.scss.tcd.ie/theses/diss/2022/TCD-SCSS-DISSERTATION-2022-016.pdf}{SPIN/Promela-Based Test Generation Framework for RTEMS Barrier Manager}] [2]
    This dissertation details previous work on the system I am currently working on. The author's contributions included improving the existing test generation framework, successfully modelling the barrier manager in RTEMS with Promela, and generating a 100\% coverage test suite for the barrier manager. This paper provided significant insight into how the system operates.

    The author identified RTEMS as a good target for formal modelling and test generation due to the fact that RTEMS as a codebase is large, having been worked on for over 30 years. as well as this, due to the fact that RTEMS was originally designed for uniprocessor systems, the test coverage for RTEMS on multi processor systems was less complete. Seeing this they saw it as a  target for the application of Promela/Spin modelling. They modelled the software in Promela and then had Spin output situations where the system fails. this is where the author steps in and created a tool that takes the output, refines it using a custom dictionary that matches specific output phrases and matches them to test code, creating a complete test file for the manager.

    Reading this I came to understand the significance of the RTEMS operating system and the importance of ensuring adequate testing. reading through it gave me the background needed to understand where the test suite generation excelled at and where it could be extended. the system was created to create tests in C. Looking at the codebase it can be seen that in many places, c specific test code is hard coded in. In view of this, my project will look to allow for greater flexibility of test code output by allowing for a language other than c to be used.

https://github.com/ddfisher/HaPy
an iniatllay very promising library that seeks to perform the exact function i needed to make this integration work, calling haskell function from python as a normal importable module. This seemed ideal until further reading of the github reveals that only basic data types are supported, so dicts and the like would not be possible to be exported with this 
// tho note to myself to look at later, do i even need to export the ref dicts, as far as i can remember they are just used as inputs, must ask the ai


https://github.com/nh2/call-haskell-from-anything
this uses an interesting system of wrapping up functions arguments and return values into MessagePacks to export via the C FFI built into haskell, seeks to improve on the base usage of the FFI by simplifying it and extending it to work for any language with an implementation of MessagePacks
\subsection{Aspect \#1}


\subsection{Aspect \#2}


 






    Connect “Why Functional Programming Matters” (Hughes) to your test generation: 

     You already have this, but you can strengthen the link:
         Hughes’ argument about modularity and higher-order functions aligns nicely with your design where:
             Lower-level functions parse YAML, handle refinement, and produce abstract test actions.
             Higher-level functions orchestrate the overall pipeline, treating those stages as black boxes.
             
         Lazy evaluation can be mentioned in the context of generating large test suites without exhausting memory (as you already do in your draft).
         




\section{Summary}

Summarize the chapter and present a comparison of the projects that you reviewed.

\begin{table}[!h]
\begin{center}
	\begin{tabular}{|l|c|c|} 
	\hline
 	\bf  & \bf Aspect \#1  & \bf Aspect \#2 \\
  	\hline
	Row 1 & Item 1 & Item 2 \\
	Row 2 & Item 1 & Item 2 \\
	Row 3 & Item 1 & Item 2 \\
	Row 4 & Item 1 & Item 2 \\
	\hline
	\end{tabular}
\end{center}
\caption[Comparison of Closely-Related Projects]{Caption that explains the table to the reader}	
\label{tab:SummaryProjects}
\end{table}
